\chapter{General Conclusion}
\label{ch:conclusion}

% ============================================================
\section{Summary of the Dissertation}
\label{sec:conclusion_summary}
% ============================================================

This dissertation investigated explainable and reproducible machine learning for \gls{ALS} prognosis through two complementary studies, both using the \gls{PRO-ACT} database and both addressing class imbalance, leakage-free validation, threshold selection, and \gls{SHAP}-based interpretability.

Study~I developed an original classification pipeline for predicting rapid versus slow \gls{ALSFRS-R} functional progression at three-month and six-month horizons. Thirty-five baseline clinical features were extracted from five \gls{PRO-ACT} data blocks and fed into seven classifiers tuned via Optuna Bayesian optimisation under 5-fold GroupKFold cross-validation with subject-level partition assignment. The pipeline was designed to be entirely leakage-free: the holdout test set was separated before any modelling, and the rapid-progressor label threshold was computed fold-wise on the training partition only. The final model, XGBoost with unbalanced class weighting, was selected based on cross-validated \gls{PR-AUC} and explained using \gls{SHAP} and \gls{LIME}.

Study~II replicated and extended the survival prognosis framework of Papaiz et al.\ \cite{Papaiz2024}, which predicts short survival (death within 24 months of symptom onset) from 23 ordinal and binary \gls{PRO-ACT} features. The replication confirmed the main claim of the reference study --- that BalancedBagging improves balanced accuracy for most classifiers --- and extended it by adding LightGBM as a seventh classifier, evaluating ten individual resampling strategies across 70 model--strategy combinations, introducing Bayesian hyperparameter optimisation with Optuna, and assessing results on a held-out test set with bootstrap confidence intervals. \gls{SHAP} attribution was applied to the two best-performing configurations using TreeSHAP and KernelSHAP respectively. A central contribution of Study~II is the empirical documentation of \gls{ROC-AUC}--sensitivity decoupling in a multi-layer perceptron under severe class imbalance, and a demonstration that this decoupling is structural rather than a threshold artefact.

% ============================================================
\section{Main Findings from Study~I}
\label{sec:conclusion_study1}
% ============================================================

The Study~I pipeline, described in Chapters~\ref{ch:study1_methods} and~\ref{ch:study1_results}, demonstrated that leakage-free prediction of \gls{ALSFRS-R} functional progression is feasible with \gls{PRO-ACT} baseline data and achieves moderate discrimination above the uninformative baseline for both the three-month and six-month horizons.

At the six-month horizon, XGBoost with unbalanced class weighting ranked first in cross-validation with a \gls{PR-AUC} of 0.434 $\pm$ 0.049, narrowly ahead of Random Forest (0.431 $\pm$ 0.042) and \gls{SVM} (0.425 $\pm$ 0.049). On the held-out test set ($n = 279$), the selected model achieved \gls{PR-AUC} = 0.456 (95\% \gls{CI}: [0.364, 0.568]). At the F1-optimal threshold ($t = 0.21$), it detected 72 of 84 rapid progressors with recall = 85.7\%, precision = 35.1\%, F1 = 0.498, and F2 = 0.665. The Brier score was 0.197, reflecting calibration that is adequate but imperfect. The three-month model performed comparably in cross-validation but showed higher variance on the held-out test, consistent with the smaller label change detectable at a shorter follow-up window.

The \gls{SHAP} analysis identified baseline \gls{ALSFRS-R} total as the primary predictor with a mean absolute \gls{SHAP} value of 0.325, approximately 2.5$\times$ larger than the second-ranked feature, \gls{FVC} in litres (0.131). Respiratory rate (0.100) and age (0.067) were the next most informative features. \gls{LIME} local explanations and logistic regression coefficients were broadly concordant with the \gls{SHAP} global ranking, which supports the clinical plausibility of the learned signal: a lower initial \gls{ALSFRS-R} score, reduced pulmonary capacity, and higher respiratory rate at baseline are all established markers of more severe disease at presentation.

The feature block ablation showed that the baseline \gls{ALSFRS-R} block alone accounts for most of the discriminative signal and that adding the \gls{FVC} block yields the most meaningful incremental gain. The remaining three blocks (vitals, treatment, and laboratory values) provided marginal contributions, which is consistent with the known correlation between respiratory function and \gls{ALSFRS-R} progression rate. Learning curves plateaued at approximately 60\% of the training data, suggesting that the performance ceiling at this sample size is feature-driven rather than data-quantity-driven. Decision curve analysis confirmed positive net benefit over the treat-all and treat-none reference strategies for a clinically plausible range of risk thresholds.

% ============================================================
\section{Main Findings from Study~II}
\label{sec:conclusion_study2}
% ============================================================

The Study~II pipeline, described in Chapters~\ref{ch:study2_methods} and~\ref{ch:study2_results}, produced three main clusters of findings: a confirmation and statistical validation of the Papaiz et al.\ \cite{Papaiz2024} BalancedBagging result; a demonstration that adding LightGBM with individual resampling and Bayesian optimisation yields better performance than the original pipeline; and an empirical characterisation of the \gls{ROC-AUC}--sensitivity decoupling phenomenon under severe class imbalance.

\textbf{BalancedBagging replication.} The Wilcoxon signed-rank test with Bonferroni correction confirmed that BalancedBagging significantly improves balanced accuracy over single-classifier configurations for six of seven tested classifiers ($p < 0.01$ after correction). The largest absolute improvements were observed for \gls{SVM} ($\Delta = +0.293$), KNN ($\Delta = +0.221$), and Random Forest ($\Delta = +0.185$). Naive Bayes was the sole non-significant case ($\Delta = +0.012$, $p > 0.05$), which is consistent with its parametric assumptions being less sensitive to the resampling applied by the bagging wrapper. These results reproduce the central claim of Papaiz et al.\ in a different software environment and with a partially different cohort, providing independent evidence for its robustness.

\textbf{Extended pipeline performance.} Among the 70 model--strategy combinations evaluated on the held-out test set ($n = 301$, 35 Short Survivors), LightGBM with random oversampling achieved the highest \gls{ROC-AUC} (0.923; 95\% \gls{CI}: [0.868, 0.966]) and a recall of 0.857 (30 of 35 short survivors detected; 95\% \gls{CI}: [0.724, 0.967]) at the default 0.5 threshold. \gls{SVM} with SMOTEENN achieved the highest recall overall (0.886, 31 of 35 detected), with \gls{ROC-AUC} = 0.909. The bootstrap confidence intervals for sensitivity overlap between LightGBM+ROS and SVM+SMOTEENN ([0.724, 0.967] versus [0.765, 0.975]), indicating that the one-patient difference is not statistically distinguishable.

\textbf{AUC--sensitivity decoupling.} The multi-layer perceptron achieved a cross-validated \gls{ROC-AUC} of 0.915 (the highest among all 70 combinations) and a test \gls{ROC-AUC} of 0.918 (95\% \gls{CI}: [0.874, 0.955]). At the default 0.5 threshold, it detected only 1 of 35 short survivors (recall = 0.029; 95\% \gls{CI}: [0.000, 0.100]). Applying Youden's J threshold (0.152) recovered recall to 0.486, but the multi-layer perceptron remained unable to match LightGBM's default-threshold recall of 0.857. The root cause is that optimising \gls{ROC-AUC} under a 7.6:1 imbalance ratio compresses minority-class posterior probabilities below 0.5 even when the model's ranking of individual patients is correct in aggregate. This finding illustrates why \gls{ROC-AUC} is insufficient as a sole evaluation criterion when the clinical objective is to detect patients from the minority class at a practical operating threshold.

\textbf{\gls{SHAP} feature attribution.} The dominant predictor in both the LightGBM and multi-layer perceptron \gls{SHAP} analyses was diagnosis delay, with a mean absolute \gls{SHAP} value of 0.986 --- more than twice the contribution of the second-ranked feature, Q8 walking slope (0.495). Longer diagnosis delay was associated with Standard Survival, whilst shorter delay and rapid decline in walking, handwriting (Q4, 0.449), and speech (Q1, 0.402) were associated with Short Survivor risk. These associations are clinically coherent: a shorter time from symptom onset to diagnosis generally reflects a more aggressive disease course.

% ============================================================
\section{Cross-Study Methodological Lessons}
\label{sec:conclusion_lessons}
% ============================================================

Several methodological conclusions emerge from examining both studies jointly, and these are likely to apply more broadly to clinical machine learning research in neurodegenerative disease.

\textbf{Target definition changes what methods are appropriate.} Study~I and Study~II address related clinical problems using the same database but different targets, features, and validation frameworks. The differences in metric choice (\gls{PR-AUC} versus \gls{ROC-AUC}), imbalance handling (class weights versus ten explicit resampling strategies), and cross-validation strategy (GroupKFold versus RepeatedStratifiedKFold) all follow from the specific characteristics of each task rather than arbitrary preference. Applying the Study~I evaluation protocol to the Study~II problem without adjustment would miss the \gls{ROC-AUC}--sensitivity decoupling, underestimate the clinical cost of missed short survivors, and report misleadingly optimistic summary statistics.

\textbf{Ranking metrics do not capture threshold-dependent performance.} Both studies provide evidence that a model with high ranking performance (\gls{ROC-AUC} or \gls{PR-AUC}) can perform poorly at a specific operating threshold. In Study~I, the XGBoost model achieves the highest cross-validated \gls{PR-AUC} but the threshold must be lowered to 0.21 (from the default 0.5) to achieve recall of 85.7\%. In Study~II, the multi-layer perceptron achieves the highest cross-validated \gls{ROC-AUC} but recalls fewer than 3\% of short survivors at the default threshold. Threshold selection using out-of-fold predictions, and subsequent reporting of recall, precision, and F1 at the selected threshold, is a necessary complement to ranking metrics in both settings.

\textbf{Leakage-free evaluation requires deliberate protocol design.} Both studies required explicit choices to prevent information leakage: subject-level holdout splits rather than record-level random splits, GroupKFold rather than simple KFold for Study~I to prevent patient records from appearing in both training and validation folds, and fold-wise computation of the rapid-progressor threshold in Study~I to prevent the label definition from encoding held-out information. These choices reduce apparent performance relative to naive splits but produce estimates that are more likely to replicate in new patient populations.

\textbf{Imbalance severity determines which corrections are sufficient.} Study~I uses class-weight adjustment, which is a low-cost intervention well suited to the 30/70 imbalance ratio. Study~II faces a 7.6:1 ratio that is severe enough for class-weight adjustment alone to be insufficient for some classifiers and for resampling strategies such as random oversampling, SMOTE, and SMOTEENN to produce substantially different results across model families. The comparison of ten resampling strategies across seven classifiers in Study~II provides a concrete demonstration that the choice of imbalance handling method is not interchangeable and interacts with classifier architecture in non-trivial ways.

\textbf{Explainability supports clinical plausibility checks, not causal inference.} In both studies, \gls{SHAP} attributions produced rankings that are coherent with clinical expectations: low baseline \gls{ALSFRS-R} and reduced \gls{FVC} as markers of rapid functional progression in Study~I, and short diagnosis delay and rapid \gls{ALSFRS-R} item decline as markers of short survival in Study~II. These results support the hypothesis that the models are learning genuine prognostic signals rather than dataset artefacts. However, feature importance from a post-hoc explainability method cannot establish that these variables causally determine outcomes, nor does it validate the model for individual patient decisions without prospective evaluation.

% ============================================================
\section{Limitations}
\label{sec:conclusion_limitations}
% ============================================================

Both studies share several limitations that bear on the interpretation of their results and on the generalisability of any conclusions drawn from them.

\textbf{Internal validation only.} Neither study includes an external validation cohort. The holdout test sets in both studies are drawn from the same \gls{PRO-ACT} extract as the training data, which means that any systematic biases in \gls{PRO-ACT} --- such as trial-protocol effects, centre heterogeneity, or selection of less severely ill patients for enrolment --- are present in both the training and test partitions. The degree to which either pipeline would generalise to clinical data from routine neurology practice or from a different national registry is unknown.

\textbf{PRO-ACT selection bias.} \gls{PRO-ACT} aggregates data from clinical trials, which tend to enrol patients with more stable disease at screening, adequate organ function, and a willingness to participate in research. This selection differs systematically from the general ALS population, and it is plausible that the performance achieved on \gls{PRO-ACT} would not replicate in an unselected outpatient cohort.

\textbf{Missing data handling.} Study~I imputes missing feature values using median substitution, a single-imputation strategy that does not propagate uncertainty from missingness into model estimates. Forced vital capacity (\gls{FVC}) is missing in approximately 47\% of records after baseline alignment, and the pattern was characterised as likely missing at random (\gls{MAR}) but not formally confirmed. Study~II applies exclusion criteria that require non-missing \gls{FVC}, which reduces the cohort by approximately 23.6\% relative to Papaiz et al.\ and introduces a selection effect towards patients with measured respiratory function.

\textbf{Small minority class in Study~II test set.} The held-out test set for Study~II contains only 35 Short Survivors. Bootstrap confidence intervals for recall are correspondingly wide: the 95\% \gls{CI} for LightGBM+ROS recall is [0.724, 0.967], spanning 24.3 percentage points. Point estimates for metrics computed on this subgroup should be interpreted with appropriate caution.

\textbf{Static baseline features.} Both studies model prognosis from a single baseline snapshot. Longitudinal feature trajectories --- repeated \gls{ALSFRS-R} measurements, changes in \gls{FVC} over the first months after diagnosis, or temporal biomarker dynamics --- were not incorporated. Longitudinal models may capture disease evolution that is not recoverable from baseline data alone.

\textbf{Absence of prospective clinical evaluation.} The per-patient risk report produced by Study~I was designed to communicate risk tier and feature attribution to clinicians, but it has not been evaluated in a clinical setting. Whether clinicians find the output actionable, whether risk tiers correspond to meaningful intervention thresholds, and whether the model's predictions are calibrated well enough to support individual decisions are open questions.

% ============================================================
\section{Future Work}
\label{sec:conclusion_future}
% ============================================================

Several directions follow directly from the limitations above and from observations made during the two studies.

\textbf{External validation.} The most immediate priority is to evaluate the Study~I pipeline on an independent ALS cohort, such as a European national registry or a real-world clinical database. This would test whether the feature block contributions, model ranking, and calibration observed on \gls{PRO-ACT} are specific to that dataset or generalisable. The Study~II pipeline could similarly be evaluated on cohorts used in other published ALS survival studies to determine whether the BalancedBagging result and the \gls{ROC-AUC}--sensitivity decoupling finding replicate.

\textbf{Multiple imputation.} Replacing median imputation in Study~I with multiple imputation by chained equations would propagate missingness uncertainty into predictions and would likely affect feature importance rankings for variables with high missingness rates. Given the extent of \gls{FVC} missingness, this is a substantive methodological concern rather than a minor refinement.

\textbf{Longitudinal modelling.} Both studies use baseline or pre-computed slope features as inputs to static classifiers. A more realistic prognostic model would update predictions as new \gls{ALSFRS-R} measurements become available, using recurrent architectures or state-space models applied to variable-length observation sequences. The \gls{PRO-ACT} database provides repeated assessments that would support this extension.

\textbf{Joint modelling of functional progression and survival.} Study~I predicts \gls{ALSFRS-R} slope and Study~II predicts mortality; the two outcomes are clinically correlated. A multi-task framework that jointly models both targets could leverage this correlation, share representations across tasks, and produce predictions that are internally consistent with respect to the expected relationship between rapid functional decline and shorter survival.

\textbf{Prospective evaluation of the risk report.} The per-patient risk report generated by the Study~I pipeline has not been evaluated by clinicians. A structured pilot study in a neurology outpatient setting would assess whether the three-tier risk classification and \gls{SHAP}-based feature attribution are interpreted correctly, whether the output format is appropriate for clinical workflow, and whether additional contextual information would be required before the tool could support individual patient decisions.

\textbf{Replication of other published pipelines.} The Study~II framework --- Bayesian optimisation, held-out test evaluation, bootstrap confidence intervals, and statistical comparison of imbalance strategies --- could be applied systematically to other published ALS prognosis pipelines. Such a programme of replication studies would contribute to establishing which methodological claims in the ALS machine learning literature are reproducible and which are specific to particular datasets or evaluation choices.

% ============================================================
\section{Final Remarks}
\label{sec:conclusion_final}
% ============================================================

Across both studies, the results show that \gls{ALS} prognostic modelling from \gls{PRO-ACT} data is achievable with standard machine learning classifiers, but that performance and clinical usefulness depend on decisions that go beyond classifier selection: how the target is defined, how temporal leakage is prevented, which metric is used to drive optimisation, how the decision threshold is selected, and what is done to explain predictions to end users.

Study~I achieved a held-out test \gls{PR-AUC} of 0.456 for six-month functional progression prediction, with a recall of 85.7\% at the F1-optimal threshold. These results are comparable to or exceed those reported in recent \gls{PRO-ACT} benchmarks \cite{Grollemund2019, Pancotti2022}, and the leakage-free design provides a more credible performance estimate than studies using naive random splits. The \gls{SHAP} analysis confirmed that baseline \gls{ALSFRS-R} total and \gls{FVC} in litres are the dominant predictors, which is consistent with their established clinical relevance.

Study~II achieved a held-out test \gls{ROC-AUC} of 0.923 for the best-performing configuration (LightGBM with random oversampling) and confirmed the statistical robustness of BalancedBagging for six of seven classifiers. The \gls{ROC-AUC}--sensitivity decoupling observed for the multi-layer perceptron at the default threshold --- test \gls{ROC-AUC} = 0.918 alongside recall = 0.029 --- is not a pathology of this particular model but a predictable consequence of optimising a ranking metric under severe class imbalance without threshold calibration. Documenting and explaining this failure mode is a contribution in itself, as it applies to any imbalanced clinical prediction problem where a minority-class label carries asymmetric clinical cost.

Neither study delivers a tool ready for clinical deployment. Both require external validation, and both rely on datasets that over-represent patients selected for clinical trial participation. What the two studies together do demonstrate is that rigorous, leakage-free, explainable, and reproducible machine learning for \gls{ALS} prognosis is tractable with publicly available data and standard open-source software, and that the methodological choices made during study design --- not the choice of classifier --- determine whether the results are interpretable, replicable, and clinically meaningful.

